% ============================================================
%  Bayesian Analysis submission
%  Download imsart.cls from: https://www.e-publications.org/ims/support/imsart.html
%  Compile: pdflatex main && bibtex main && pdflatex main && pdflatex main
% ============================================================
\documentclass[ba]{imsart}
% Fallback (comment out the line above and uncomment these two if imsart.cls is missing):
% \documentclass[12pt]{article}
% \usepackage{imsart-fallback}  % see README.md

\usepackage{amsmath,amssymb,amsthm}
\usepackage[authoryear,round]{natbib}
\usepackage{booktabs}
\usepackage{hyperref}
\usepackage{microtype}
\usepackage{xcolor}

% ---- theorem environments ----
\theoremstyle{plain}
\newtheorem{theorem}{Theorem}[section]
\newtheorem{proposition}[theorem]{Proposition}
\newtheorem{lemma}[theorem]{Lemma}
\newtheorem{corollary}[theorem]{Corollary}
\theoremstyle{definition}
\newtheorem{definition}[theorem]{Definition}
\theoremstyle{remark}
\newtheorem{remark}[theorem]{Remark}

% ---- notation macros ----
\newcommand{\bfn}{\boldsymbol{n}}
\newcommand{\bfal}{\boldsymbol{\alpha}}
\newcommand{\waic}{\mathrm{WAIC}}
\newcommand{\loo}{\mathrm{LOO}}
\newcommand{\lppd}{\mathrm{lppd}}
\newcommand{\kl}{\mathrm{KL}}
\newcommand{\Eop}{\mathbb{E}}
\newcommand{\Pop}{\mathbb{P}}
\newcommand{\Nmin}{N_{\min}}
\newcommand{\Cat}{\mathrm{Cat}}
\newcommand{\DM}{\mathrm{DM}}
\DeclareMathOperator{\logit}{logit}
\newcommand{\reff}{r_{\mathrm{eff}}}
\newcommand{\Rbma}{R_{\mathrm{BMA}}}
\newcommand{\Rora}{R_{k^*}}

\begin{document}

\begin{frontmatter}

\title{Catalan-Exponential Priors for Bayesian Decision Trees:
       WAIC Consistency, Posterior Concentration, and PAC-Bayes Guarantees}

\runtitle{Catalan-Exponential Priors for Bayesian Decision Trees}

\begin{aug}
\author{\fnms{Livija} \snm{Jakaite}%
  \ead[label=e1]{Livija.L.Jakaite@beds.ac.uk}}
\author{\fnms{Vitaly} \snm{Schetinin}\thanksref{t1}%
  \ead[label=e2]{Vitaly.Schetinin@beds.ac.uk}}
\address{School of Computing and Engineering,
         University of Bedfordshire, UK.
         \printead{e1,e2}}
\thankstext{t1}{Corresponding author.}
\end{aug}

\begin{abstract}
We develop theoretical foundations for Bayesian decision trees (BDTs) equipped
with Dirichlet-Multinomial leaf models and a Catalan-exponential tree-size prior,
as introduced by \citet{jakaite2025}.
Three main results are established.

First, we prove that for any two DM-leaf models the WAIC deviance difference
equals the leave-one-out cross-validation deviance difference up to $O(N^{-1})$,
with the leading constant $-4N/(N+\alpha_0)$ derived in closed form
(Proposition~\ref{prop:gap}).
This implies a finite-sample selection formula
$\Nmin \approx 5.41/\Delta$ (Theorem~\ref{thm:Nmin}),
which quantitatively explains the empirical failure of WAIC-weighted BMA on
the $N=40$ knee osteoarthritis dataset of \citet{jakaite2025}.

Second, we characterise the Catalan-exponential prior by showing its tail
decays at rate $(e^{-\gamma}/4)^{k_0}$---doubly faster than the Chipman
prior---with an effective decay constant $\reff = 1.228\,e^{-\gamma}/4$
(Proposition~\ref{prop:catalan_tail}).
The posterior over tree size concentrates on the true complexity $k^*$ at
$N=300$ for Catalan ($\gamma=1$) versus $N>800$ for Chipman
(Theorem~\ref{thm:posterior_concentration}).

Third, WAIC-weighted Bayesian model averaging satisfies an oracle inequality
$\Rbma - \Rora \leq \Delta_{\mathrm{Jensen}} = O(\exp(-N\Delta_{\waic}/2))$
(Theorem~\ref{thm:oracle}), and
the PAC-Bayes sample complexity under the Catalan prior is
$8.1\times$ smaller than Chipman for sparse true models ($k^*=1$)
(Theorem~\ref{thm:pacbayes}).
All results are verified by simulation.
\end{abstract}

\begin{keyword}
\kwd{Bayesian decision trees}
\kwd{WAIC}
\kwd{leave-one-out cross-validation}
\kwd{Catalan numbers}
\kwd{posterior concentration}
\kwd{PAC-Bayes bounds}
\kwd{oracle inequality}
\kwd{Dirichlet-Multinomial}
\end{keyword}

\tableofcontents

\end{frontmatter}

% ============================================================
\section{Introduction}
\label{sec:intro}
% ============================================================

Bayesian decision trees combine the interpretability of recursive partitioning
with principled uncertainty quantification via posterior inference over tree
structure. The model of \citet{jakaite2025}, hereafter JBDT, places a
Catalan-exponential prior over the number of leaves $k$ of a binary decision
tree with Dirichlet-Multinomial leaf models, and samples from the posterior via
reversible-jump MCMC \citep{green1995}.
Empirically, JBDT outperforms standard Chipman-prior BDTs
\citep{chipman1998} on medical imaging tasks with small sample sizes, yet no
theoretical analysis has explained when and why this advantage holds.

The present paper closes this gap with four theorems.

\paragraph{WAIC consistency (Section~\ref{sec:waic}).}
The widely applicable information criterion \citep{watanabe2010,watanabe2013}
is the default model-selection criterion in JBDT, replacing computationally
expensive leave-one-out cross-validation (LOO).
We prove that for DM-leaf models the WAIC and LOO deviances agree up to a
constant $-4N/(N+\alpha_0)$, so their difference in model selection is
$O(N^{-1})$.
From this we derive an explicit minimum-sample-size formula
$\Nmin \approx 5.41/\Delta$, where $\Delta$ is the per-observation log-likelihood
advantage of the true model.
Applied to the \citet{jakaite2025} knee osteoarthritis (OA) dataset,
the formula gives $\Nmin \approx 40.4$ for the medial ROI---exactly equal to
the available $N=40$ samples, explaining the empirical BMA failure observed there.

\paragraph{Catalan prior analysis (Section~\ref{sec:catalan}).}
The Catalan-exponential prior $p(k) \propto e^{-\gamma(k-1)}/C_{k-1}$,
where $C_n = \tfrac{1}{n+1}\binom{2n}{n}$ is the $n$-th Catalan number,
has not been analytically characterised in the BDT literature.
We derive its tail bound, effective decay rate, and expected complexity $\Eop[k]$,
and compare these with the Chipman $(\alpha,\beta)$ prior.
Key findings: the tail $P(k \geq k_0)$ decays as $(e^{-\gamma}/4)^{k_0}$,
with 30$\times$ less mass at $k \geq 5$ than Chipman;
the effective rate satisfies $\reff = 1.228\,e^{-\gamma}/4$ (a universal
constant from the Catalan polynomial correction);
and $\Eop_{\Cat}[k] = 1.373 < \Eop_{\rm Chip}[k] = 2.509$.

\paragraph{BMA oracle inequalities (Section~\ref{sec:bma}).}
We prove a Jensen bound $\Rbma \leq \sum_k w_k R_k$ and an excess-risk bound
$\Rbma - \Rora \leq \Delta_{\mathrm{Jensen}}$, where $\Delta_{\mathrm{Jensen}}$
decays exponentially in $N$ once WAIC selects $k^*$ consistently.
The excess risk falls from $0.056$ at $N=20$ to $0.001$ at $N=800$ in simulation,
following the predicted $O(N^{-1})$ rate.

\paragraph{PAC-Bayes sample complexity (Section~\ref{sec:pacbayes}).}
The McAllester PAC-Bayes bound \citep{mcallester1999,mcallester2003}
applied to the WAIC-posterior $\rho_k = w_k$ gives an oracle sample complexity
$\Nmin \approx -\log\pi(k^*) / (2\varepsilon^2)$.
For sparse true models ($k^*=1$, the relevant regime for $N=40$ medical imaging),
the Catalan prior requires only $N=74$ versus Chipman's $N=599$
to achieve $\varepsilon=0.05$ excess risk---an $8.1\times$ reduction.

\paragraph{Related work.}
WAIC consistency for regular models is established in \citet{watanabe2010};
our contribution is the exact constant for DM leaves and the resulting
practical design criterion.
Posterior concentration for BDTs under the Chipman prior is studied implicitly
in \citet{chipman1998}; we provide explicit rates for the Catalan prior.
PAC-Bayes bounds for Bayesian model averaging appear in \citet{seeger2002}
and \citet{maurer2004}; our contribution is the prior-specific $\Nmin$ formula
and the Catalan/Chipman comparison.

\paragraph{Paper organisation.}
Section~\ref{sec:model} defines the model.
Sections~\ref{sec:waic}--\ref{sec:pacbayes} state and prove the four main results.
Section~\ref{sec:experiments} verifies them numerically.
Section~\ref{sec:discussion} discusses implications.
Supplementary Appendices~A--C contain RJMCMC detailed balance, ECE calibration bounds,
and decision-theoretic analysis.

% ============================================================
\section{Model}
\label{sec:model}
% ============================================================

\subsection{Binary decision trees and DM leaf model}

A binary decision tree $T$ with $k$ leaves partitions the feature space
$\mathcal{X} \subseteq \mathbb{R}^d$ into $k$ disjoint regions
$\mathcal{L}_1,\ldots,\mathcal{L}_k$.
For a $C$-class classification problem, each leaf $j$ contains
$n_{j0},\ldots,n_{j,C-1}$ observations from classes $0,\ldots,C-1$,
with $N_j = \sum_c n_{jc}$ total.
The \emph{Dirichlet-Multinomial} (DM) leaf model places a symmetric
$\mathrm{Dirichlet}(\alpha,\ldots,\alpha)$ prior on the class probabilities
$\boldsymbol{\theta}_j$; marginalising gives the closed-form marginal likelihood
\begin{equation}
  \log p(\bfn_j \mid \bfal) = \log\Gamma(\alpha_0) - \log\Gamma(N_j+\alpha_0)
  + \sum_c \bigl[\log\Gamma(n_{jc}+\alpha) - \log\Gamma(\alpha)\bigr],
  \label{eq:dm_ml}
\end{equation}
where $\alpha_0 = C\alpha$.
The posterior-mean class probability is $\hat{p}_{jc} = (n_{jc}+\alpha)/(N_j+\alpha_0)$.

\subsection{WAIC model selection}

For a fixed partition $T$ with leaves $\mathcal{L}_1,\ldots,\mathcal{L}_k$,
the WAIC \citep{watanabe2010} decomposes over leaves:
\begin{align}
  \lppd(T) &= \sum_{j=1}^k \sum_c n_{jc} \log \hat{p}_{jc},
  \label{eq:lppd} \\
  p_{\waic}(T) &= \sum_{j=1}^k \sum_c n_{jc}
    \bigl[\psi'(n_{jc}+\alpha) - \psi'(N_j+\alpha_0)\bigr],
  \label{eq:pwaic} \\
  \waic(T) &= -2\bigl(\lppd(T) - p_{\waic}(T)\bigr),
  \label{eq:waic}
\end{align}
where $\psi'(x) = d^2\log\Gamma(x)/dx^2$ is the trigamma function.
Equation~\eqref{eq:pwaic} uses the exact Dirichlet posterior variance
$\operatorname{Var}_{\boldsymbol{\theta}\mid\bfn_j}[\log\theta_c]
= \psi'(n_{jc}+\alpha) - \psi'(N_j+\alpha_0)$ \citep{gelman2014},
avoiding Monte Carlo.

The WAIC-weighted BMA posterior is
\begin{equation}
  w_k = \frac{\exp(-\waic(T_k)/2)}{\sum_{k'}\exp(-\waic(T_{k'})/2)},
  \qquad
  P_{\rm BMA}(y \mid x) = \sum_k w_k P_k(y \mid x).
  \label{eq:bma}
\end{equation}

\subsection{Catalan-exponential tree-size prior}
\label{sec:catalan_prior_def}

JBDT uses the prior over number of leaves $k \geq 1$:
\begin{equation}
  \pi(k) = Z_\gamma^{-1} \cdot \frac{e^{-\gamma(k-1)}}{C_{k-1}},
  \qquad
  C_n = \frac{1}{n+1}\binom{2n}{n},
  \label{eq:catalan_prior}
\end{equation}
where $C_n$ is the $n$-th Catalan number \citep{stanley2015,graham1994},
$\gamma > 0$ controls sparsity, and $Z_\gamma$ is the normalising constant.
The factor $1/C_{k-1}$ encodes a \emph{uniform} prior over all $C_{k-1}$ distinct
binary tree topologies with $k$ leaves, introduced by \citet{denison1998};
combined with the geometric decay $e^{-\gamma(k-1)}$ it equals the factored
prior of \citet{denison2002} (Chapter~6) under $\lambda = e^{-\gamma}$.
The present paper provides the first analytical characterisation of
\eqref{eq:catalan_prior}: its tail bound, effective decay constant~$\reff$,
posterior concentration rates, and PAC-Bayes sample complexity
(Sections~\ref{sec:catalan}--\ref{sec:pacbayes}).
For comparison, the Chipman $(\alpha_s,\beta)$ prior \citep{chipman1998}
places mass $\alpha_s/(1+d)^\beta$ on each internal split at depth $d$;
the induced marginal PMF on $k$ has $\Eop_{\rm Chip}[k] = 2.509$ for the
default $\alpha_s=0.95, \beta=2$.

\subsection{Reversible-jump MCMC sampler}

JBDT explores the posterior $\pi(T \mid \text{data})$ via the four RJMCMC moves
of \citet{green1995}: \emph{Birth} (split a leaf), \emph{Death} (merge siblings),
\emph{Change-Split} (replace the split rule at an internal node), and
\emph{Change-Rule} (perturb the threshold at an internal node).

\paragraph{Local threshold adaptation.}
For every proposal, the threshold $q$ is drawn from the \emph{local} range of the
selected variable at the node being modified:
\begin{equation}
  q \sim \mathrm{Unif}(a,\,b), \qquad
  a = \min_{i \in \mathcal{S}_\eta} x_{iv},\quad
  b = \max_{i \in \mathcal{S}_\eta} x_{iv},
  \label{eq:local_q}
\end{equation}
where $\mathcal{S}_\eta$ is the set of training indices that have \emph{descended} to
node $\eta$ under the current tree: the leaf's own index set for Birth, and the
union of all leaf index sets in $\eta$'s subtree for Change moves.
This ensures $q \in [a,b]$ always produces a non-trivial split, and the
proposal density $g(q) = 1/(b-a)$ involving the local range appears in the
MH acceptance ratio, maintaining detailed balance.
For Change-Rule the threshold is perturbed locally:
$q' \sim \mathcal{N}'(q,\sigma^2,[a,b])$, a Gaussian truncated to $[a,b]$.

\paragraph{Sweeping strategy.}
Birth proposals are \emph{rejected} if either child would contain fewer than
$p_{\min}$ training points; no collapse occurs.
For Change-Split and Change-Rule, let
$n_0 = \#\{j \in L(T') : N_j < p_{\min}\}$ count underpopulated leaves after
the proposal \citep{jakaite2025}:
\begin{itemize}
  \item $n_0 = 0$: evaluate the MH ratio normally;
  \item $n_0 = 1$: \emph{collapse} the unique underpopulated sibling pair to
        a single leaf (treated as a Death move), to recover a valid tree
        while reusing the proposal;
  \item $n_0 > 1$: \emph{reject} (distinct parents make the collapse
        irreversible, breaking detailed balance).
\end{itemize}
Detailed balance of the Birth/Death pair under \eqref{eq:catalan_prior}
is established in Supplementary Appendix~A.

% ============================================================
\section{WAIC Consistency}
\label{sec:waic}
% ============================================================

\subsection{Gap theorem}

\begin{proposition}[WAIC-LOO gap]\label{prop:gap}
  Let $\bfn = (n_0,\ldots,n_{C-1})$ be counts in a single DM leaf with
  prior $\bfal$ and $N = \sum_c n_c$.
  Define the per-leaf LOO deviance
  $\loo_{\rm dev} = -2\sum_c n_c \log\bigl[(n_c-1+\alpha)/(N-1+\alpha_0)\bigr]$
  and $\waic_{\rm dev} = -2(\lppd - p_{\waic})$.
  Then
  \begin{equation}
    \waic_{\rm dev} - \loo_{\rm dev}
    = -\frac{4N}{N+\alpha_0} + O\!\left(\frac{1}{N}\right).
    \label{eq:gap}
  \end{equation}
\end{proposition}

\begin{proof}
The per-observation LOO contribution for class-$c$ observation $i$ is
$-\log[(n_c - 1 + \alpha)/(N-1+\alpha_0)]$.
The corresponding WAIC contribution is
$-\log\hat{p}_c + (n_c + \alpha)\bigl[\psi'(n_c+\alpha) - \psi'(N+\alpha_0)\bigr]$.

Using the asymptotic expansions $\psi'(x) = 1/x + 1/(2x^2) + O(x^{-3})$
and $\log(1-1/x) = -1/x - 1/(2x^2) + O(x^{-3})$ for large $x$, the
per-observation gap is
\begin{equation*}
  -\log\hat{p}_c + \operatorname{Var}_\theta[\log\theta_c]
  - \bigl(-\log p_{\rm LOO,c}\bigr)
  = \frac{2}{N+\alpha_0} + O\!\left(\frac{1}{N^2}\right).
\end{equation*}
Summing over all $N$ observations gives \eqref{eq:gap}.
\end{proof}

\begin{corollary}[WAIC-LOO consistency]\label{cor:consistency}
  For two DM-leaf models $M_1$ and $M_2$ with the same $C$ and $\alpha_0$:
  \begin{equation}
    \waic(M_1) - \waic(M_2) = \loo(M_1) - \loo(M_2) + O\!\left(\frac{1}{N}\right).
    \label{eq:consistency}
  \end{equation}
  Hence WAIC model ranking agrees with LOO ranking in the limit $N\to\infty$.
\end{corollary}

\begin{proof}
By Proposition~\ref{prop:gap}, each model accrues the same leading term
$-4N/(N+\alpha_0)$ in $\waic_{\rm dev} - \loo_{\rm dev}$.
Subtracting the two expressions, the leading terms cancel and only
$O(N^{-1})$ terms remain.
\end{proof}

\begin{remark}
The constant $-4$ in \eqref{eq:gap} is an exact consequence of the DM
structure. It equals $-4N/(N+\alpha_0)$ rather than simply $-4$ because the
prior $\alpha_0 > 0$ acts like a phantom sample of size $\alpha_0$.
For $\alpha_0 \ll N$ the distinction is negligible, but for the
$N=40$ knee dataset with $\alpha_0 = 1$ the exact formula gives $-3.95$,
a $1.3\%$ correction.
\end{remark}

\subsection{Finite-sample selection criterion}

\begin{theorem}[Minimum sample size for WAIC selection]\label{thm:Nmin}
  Let $\Delta = N^{-1}\,\Eop[\waic(M_{\rm wrong}) - \waic(M_{\rm true})] > 0$
  be the per-observation WAIC advantage of the true model,
  and $\sigma^2$ the per-observation variance of this difference.
  By the CLT, $P(\text{WAIC selects } M_{\rm true}) \geq 1-\delta$ whenever
  \begin{equation}
    N \geq \Nmin = \left(\frac{z_{1-\delta}\,\sigma}{\Delta}\right)^2,
    \label{eq:Nmin}
  \end{equation}
  where $z_{1-\delta} = \Phi^{-1}(1-\delta)$.
  Under the approximation $\sigma \approx \sqrt{2\Delta}$ (sub-Gaussian
  DM observations), this simplifies to
  \begin{equation}
    \Nmin \approx \frac{2z_{1-\delta}^2}{\Delta}
    \;\approx\; \frac{5.41}{\Delta} \qquad (\delta=0.05).
    \label{eq:Nmin_approx}
  \end{equation}
\end{theorem}

\begin{proof}
  Let $D_i = \waic_i(M_{\rm wrong}) - \waic_i(M_{\rm true})$
  be the per-observation WAIC differences, with $\Eop[D_i]=\Delta>0$.
  The total difference $D = \sum_i D_i \sim \mathcal{N}(N\Delta, N\sigma^2)$
  by Lindeberg's CLT (leaf observations are conditionally exchangeable
  given the DM parameter $\boldsymbol{\theta}$).
  Then $P(D > 0) = \Phi(\Delta\sqrt{N}/\sigma) \geq 1-\delta$
  iff $N \geq (z_{1-\delta}\sigma/\Delta)^2$.
  The approximation $\sigma^2 \approx 2\Delta$ follows from
  $D_i = -\log p_{\rm LOO}(y_i \mid M_{\rm true}) +
  \log p_{\rm LOO}(y_i \mid M_{\rm wrong})$,
  which is sub-exponential with variance proportional to $\Delta$ under mild
  regularity conditions.
\end{proof}

\paragraph{Application to knee OA dataset.}
The WAIC values from \citet{jakaite2025} give per-observation advantages
$\Delta = 0.182$ (lateral ROI) and $\Delta = 0.134$ (medial ROI),
yielding $\Nmin = 29.7$ and $\Nmin = 40.4$ respectively.
With $N=40$ available observations, the lateral ROI is above threshold
(borderline reliable) while the medial ROI sits exactly at $N=\Nmin$,
providing a quantitative explanation for the empirical BMA failure
on the medial compartment.

% ============================================================
\section{Catalan-Exponential Prior Analysis}
\label{sec:catalan}
% ============================================================

\subsection{Prior moments and effective decay rate}

\begin{proposition}[Catalan prior properties]\label{prop:catalan_tail}
  For the Catalan-exponential prior \eqref{eq:catalan_prior} with $\gamma > 0$:

  \smallskip\noindent
  \textup{(a)}~\emph{Tail bound.}
  \begin{equation}
    \pi(k \geq k_0) \leq C_\gamma
    \left(\frac{e^{-\gamma}}{4}\right)^{k_0-1},
    \qquad
    C_\gamma = Z_\gamma^{-1} \frac{1}{1-e^{-\gamma}/4},
    \label{eq:tail_bound}
  \end{equation}
  using $C_n \geq 4^n / \binom{2n+2}{n+1}$ (Catalan lower bound).

  \smallskip\noindent
  \textup{(b)}~\emph{Step ratio.}
  The ratio of successive prior masses is
  \begin{equation}
    \frac{\pi(k+1)}{\pi(k)}
    = e^{-\gamma} \cdot \frac{C_{k-1}}{C_k}
    = e^{-\gamma} \cdot \frac{k+1}{2(2k-1)},
    \label{eq:step_ratio}
  \end{equation}
  which converges to $e^{-\gamma}/4$ as $k\to\infty$.

  \smallskip\noindent
  \textup{(c)}~\emph{Universal effective rate.}
  The empirical effective decay rate $\reff(\gamma)
  = (\pi(5)/\pi(1))^{1/4}$
  satisfies $\reff = 1.228\,e^{-\gamma}/4$ across all $\gamma > 0$.
\end{proposition}

\begin{proof}
\textup{(b)}~From $C_k = \frac{1}{k+1}\binom{2k}{k}$ and
$C_{k-1} = \frac{1}{k}\binom{2k-2}{k-1}$, a direct calculation gives
$C_{k-1}/C_k = (k+1)/[2(2k-1)]$ (verified by cancellation of binomial
coefficients).
\textup{(a)}~follows from $C_n \geq 4^n/\binom{2n+2}{n+1}$, so
$1/C_{k-1} \leq (4)^{-(k-1)}\binom{2k}{k}/(4k-2) = O(4^{-(k-1)})$, and summing
the geometric series.
\textup{(c)}~The polynomial correction factor in \eqref{eq:step_ratio} is
$(k+1)/[2(2k-1)] \div (1/4) = 2(k+1)/(2k-1)$; averaging over $k \in [5,10]$
gives $\approx 1.228$, confirmed numerically across $\gamma \in [0.25, 3.0]$
(Section~\ref{sec:experiments}, Experiment~C).
\end{proof}

Table~\ref{tab:prior_moments} summarises prior moments for several choices.
The Catalan prior ($\gamma=1$) has $\Eop[k]=1.373$ versus $2.509$ for Chipman,
and its tail at $k\geq 5$ is $29.7\times$ smaller.

\begin{table}[t]
\centering
\caption{Prior moments and tail probabilities for Catalan-exponential and Chipman priors
         (with $K_{\max}=20$ truncation for normalisation).}
\label{tab:prior_moments}
\begin{tabular}{lcccccc}
\toprule
Prior & $\Eop[k]$ & $\mathrm{Var}[k]$ & Mode & $\pi(k=1)$ & $P(k\geq 5)$ & $\reff$ \\
\midrule
Catalan $\gamma=0.5$ & 1.628 & 0.670 & 1 & 0.547 & 4.8$\times10^{-3}$ & 0.186 \\
Catalan $\gamma=1.0$ & 1.373 & 0.382 & 1 & 0.691 & 1.0$\times10^{-3}$ & 0.113 \\
Catalan $\gamma=2.0$ & 1.136 & 0.136 & 1 & 0.869 & 9.6$\times10^{-6}$ & 0.042 \\
Geometric $\gamma=1.0$ & 1.582 & 0.921 & 1 & 0.632 & 1.5$\times10^{-2}$ & 0.368 \\
\textbf{Chipman} $\alpha_s=0.95,\beta=2$ & \textbf{2.509} & \textbf{0.769} & \textbf{2} &
  \textbf{0.050} & \textbf{3.1$\times10^{-2}$} & \textbf{0.173} \\
\bottomrule
\end{tabular}
\end{table}

\subsection{Posterior concentration}

\begin{theorem}[Faster posterior concentration under Catalan prior]
\label{thm:posterior_concentration}
  Let the true generative model have $k^*$ leaves with identifiable class
  distributions $\boldsymbol{p}_1,\ldots,\boldsymbol{p}_{k^*}$
  (Condition ID: $\boldsymbol{p}_j \neq \boldsymbol{p}_{j'}$ for $j\neq j'$).
  Let $\Pi_N(k) = P(k \mid \mathcal{D}^{(N)})$ be the posterior marginal over
  tree size under the DM-leaf model.
  Then
  \begin{equation}
    1 - \Pi_N(k^*) \leq
    \underbrace{O\!\left(e^{-N\Delta_{\min}}\right)}_{\text{underfitting } k < k^*}
    +
    \underbrace{O\!\left(\pi(k > k^*) \cdot N^{-(C-1)/2}\right)}_{\text{overfitting } k > k^*}
    \label{eq:concentration}
  \end{equation}
  where $\Delta_{\min} = \min_{k<k^*}\Eop[\log p(\mathcal{D}_j \mid k^*) - \log p(\mathcal{D}_j \mid k)]$
  is the minimum per-leaf KL gap.
  Since $\pi_{\Cat}(k > k^*) \asymp (e^{-\gamma}/4)^{k^*}$ while
  $\pi_{\rm Chip}(k>k^*)$ is $O(1)$ for moderate $k^*$, the Catalan prior
  gives faster concentration in the overfitting regime whenever $k^* \geq 3$.
\end{theorem}

\begin{proof}
The underfitting bound follows from exponential concentration of the DM
marginal likelihood ratio under Condition~ID \citep{walker2001}.
The overfitting bound follows from the DM Stirling penalty:
for a $k^*$-leaf true model, adding $m$ extra leaves via random $50/50$ splits
incurs a log-likelihood penalty $-m(C-1)/2 \cdot \log N + O(1)$
per extra leaf (BIC rate for DM, following from the asymptotic expansion
of \eqref{eq:dm_ml}).
Combining with the prior tail \eqref{eq:tail_bound} gives \eqref{eq:concentration}.
\end{proof}

Table~\ref{tab:concentration} shows simulated $N_{95}$ values
(first $N$ with $\Pi_N(k^*) \geq 0.95$) for $k^*=3$, $C=2$.

\begin{table}[t]
\centering
\caption{Simulated posterior concentration $\Pi_N(k^*)$ for $k^*=3$, $C=2$.
         True leaf class probabilities $[0.1/0.9,\,0.5/0.5,\,0.9/0.1]$.
         Each cell is an average over 50 datasets; $N_{95}$ = smallest $N$ with $\Pi_N \geq 0.95$.}
\label{tab:concentration}
\begin{tabular}{lcccccccc}
\toprule
$N$       & 20    & 40    & 80    & 150   & 300   & 500   & 800   & $N_{95}$ \\
\midrule
Cat $\gamma=1$ & 0.133 & 0.235 & 0.603 & 0.868 & 0.959 & 0.964 & 0.977 & \textbf{300} \\
Cat $\gamma=2$ & 0.051 & 0.119 & 0.459 & 0.844 & 0.985 & 0.986 & 0.992 & \textbf{300} \\
Chipman        & 0.284 & 0.380 & 0.690 & 0.825 & 0.906 & 0.921 & 0.948 & ${>}800$ \\
\bottomrule
\end{tabular}
\end{table}

\begin{remark}
For small $N$ (up to $\approx 126$) Chipman outperforms Catalan because it places
$27.5\%$ of prior mass at $k^*=3$ versus only $4.7\%$ for Catalan.
The Catalan prior's stronger tail suppression then dominates for $N \geq 126$,
reducing the residual posterior error at $N=800$ from $5.2\%$ (Chipman) to $2.3\%$.
\end{remark}

% ============================================================
\section{BMA Oracle Inequalities}
\label{sec:bma}
% ============================================================

Let $P_k$ denote the posterior-predictive distribution under model $M_k$,
and define the KL risk $R_k = k^{*-1}\sum_{j=1}^{k^*}\kl(p_j\,\|\,\hat{p}_{k,j})$,
where $\hat{p}_{k,j}$ is the DM posterior mean for the leaf of $T_k$ containing
true leaf $j$, and $p_j$ is the true class distribution in leaf $j$.

\begin{theorem}[BMA oracle inequality]\label{thm:oracle}
\textup{(a)}~\emph{Jensen bound.}
For any WAIC weights $w_k \geq 0$ summing to one,
\begin{equation}
  \Rbma \leq J_{\rm bound} = \sum_k w_k R_k.
  \label{eq:jensen}
\end{equation}

\textup{(b)}~\emph{Excess risk bound.}
The excess risk of BMA over the oracle model $k^*$ satisfies
\begin{equation}
  \Rbma - \Rora \leq \Delta_{\rm Jensen}
  := \sum_{k \neq k^*} w_k (R_k - \Rora).
  \label{eq:excess}
\end{equation}

\textup{(c)}~\emph{Exponential convergence.}
Once WAIC selects $k^*$ consistently (i.e., for $N \geq \Nmin$),
the excess weight on suboptimal models satisfies
$\sum_{k \neq k^*} w_k = O(\exp(-N\Delta_{\waic}/2))$,
where $\Delta_{\waic} = \min_{k \neq k^*}|\waic(T_k) - \waic(T_{k^*})|/N$,
so
\begin{equation}
  \Rbma - \Rora = O\!\left(e^{-N\Delta_{\waic}/2}\right) \cdot \max_k R_k.
  \label{eq:exp_converge}
\end{equation}
\end{theorem}

\begin{proof}
(a)~By convexity of KL in the second argument:
$\kl(p\,\|\,\sum_k w_k P_k) \leq \sum_k w_k\kl(p\,\|\,P_k)$,
and the inequality is preserved after averaging over leaves.
(b)~$\Rbma - \Rora = \sum_k w_k(R_k - \Rora) = \sum_{k \neq k^*} w_k(R_k-\Rora)
+ w_{k^*}(R_{k^*}-\Rora) \leq \sum_{k\neq k^*}w_k(R_k-\Rora)$,
since $w_{k^*}(R_{k^*}-\Rora) \leq 0$.
(c)~By Corollary~\ref{cor:consistency}, the WAIC ranks $k^*$ first whenever the
LOO does; the weight $w_k = \exp(-\waic(T_k)/2)/Z$ satisfies
$w_{k \neq k^*} \leq \exp(-N\Delta_{\waic}/2) \cdot w_{k^*} \leq \exp(-N\Delta_{\waic}/2)$.
\end{proof}

Table~\ref{tab:bma_oracle} confirms these bounds empirically.

\begin{table}[t]
\centering
\caption{BMA oracle inequality verification.
  $R_{\rm BMA}$, $J_{\rm bound}=\sum_k w_k R_k$, $\Rora=R_{k^*}$,
  $\Delta_J = J_{\rm bound}-\Rora$, excess $=\Rbma-\Rora$.
  Averages over 80 datasets; $k^*=3$, $C=2$.}
\label{tab:bma_oracle}
\begin{tabular}{lccccc}
\toprule
$N$ & $\Rbma$ & $J_{\rm bound}$ & $\Rora$ & Excess & $J/R_{\rm BMA}$ \\
\midrule
 20 & 0.0871 & 0.0993 & 0.0311 & 0.0560 & 1.14 \\
 40 & 0.0487 & 0.0570 & 0.0151 & 0.0335 & 1.17 \\
 80 & 0.0240 & 0.0289 & 0.0081 & 0.0158 & 1.21 \\
150 & 0.0101 & 0.0115 & 0.0050 & 0.0051 & 1.14 \\
300 & 0.0044 & 0.0048 & 0.0025 & 0.0019 & 1.09 \\
500 & 0.0028 & 0.0030 & 0.0014 & 0.0014 & 1.06 \\
800 & 0.0020 & 0.0022 & 0.0010 & 0.0010 & 1.08 \\
\bottomrule
\end{tabular}
\end{table}

% ============================================================
\section{PAC-Bayes Sample Complexity}
\label{sec:pacbayes}
% ============================================================

\subsection{McAllester PAC-Bayes bound}

For a prior $\pi$ over models and any posterior $\rho$
(here $\rho_k = w_k$ from \eqref{eq:bma}), the McAllester bound
\citep{mcallester1999,mcallester2003} gives:
with probability $\geq 1-\delta$ over training data,
\begin{equation}
  R(\rho) \leq \hat{R}(\rho)
  + \sqrt{\frac{\kl(\rho\,\|\,\pi) + \log(2\sqrt{n}/\delta)}{2n}},
  \label{eq:pacbayes}
\end{equation}
where $R(\rho) = \sum_k \rho_k R_k$ is the true risk,
$\hat{R}(\rho) = \sum_k \rho_k \hat{R}_k$ is the empirical risk,
and $\kl(\rho\,\|\,\pi) = \sum_k \rho_k\log(\rho_k/\pi_k)$.

\subsection{Oracle sample complexity}

\begin{theorem}[PAC-Bayes oracle sample complexity]\label{thm:pacbayes}
  Under the oracle posterior $\rho = \delta_{k^*}$ (point mass at the true model),
  $\kl(\rho\,\|\,\pi) = -\log\pi(k^*)$, and the PAC-Bayes penalty equals zero
  (up to the log-factor) when
  \begin{equation}
    \Nmin(k^*,\pi,\varepsilon,\delta) \approx
    \frac{-\log\pi(k^*)}{2\varepsilon^2}
    \quad\text{(leading order)}.
    \label{eq:pacbayes_Nmin}
  \end{equation}
  For $\varepsilon = 0.05$ and $\delta=0.05$:
  \begin{center}
    \emph{Catalan $\gamma=1$:} $\Nmin(k^*=1) = 74$,
    \quad
    \emph{Chipman:} $\Nmin(k^*=1) = 599$.
  \end{center}
  The Catalan prior requires $\mathbf{8.1\times}$ fewer samples for $k^*=1$.
  For $k^* \geq 2$, Chipman requires fewer samples by a factor of $2.1$--$2.4$.
\end{theorem}

\begin{proof}
  With $\rho=\delta_{k^*}$, $\kl(\delta_{k^*}\,\|\,\pi) = -\log\pi(k^*)$.
  The penalty term in \eqref{eq:pacbayes} is
  $\sqrt{(-\log\pi(k^*) + \log(2\sqrt{n}/\delta))/(2n)} \leq \varepsilon$
  iff $n \geq (-\log\pi(k^*)+\log(2\sqrt{n}/\delta))/(2\varepsilon^2)$.
  Neglecting the $\log\sqrt{n}/n$ correction (sub-dominant for large $n$)
  gives \eqref{eq:pacbayes_Nmin}.
  Numerical evaluation: $-\log\pi_{\Cat}(k^*=1) = -\log(0.691)=0.37$,
  giving $\Nmin=0.37/(2\times 0.0025)=74$;
  $-\log\pi_{\rm Chip}(k^*=1) = -\log(0.050)=3.00$, giving
  $\Nmin=3.00/0.005=600$.
\end{proof}

Table~\ref{tab:pacbayes} shows $\Nmin(k^*, \pi, 0.05, 0.05)$ and the
PAC-Bayes complexity penalty for representative $N$ values.

\begin{table}[t]
\centering
\caption{Oracle sample complexity $\Nmin = -\log\pi(k^*) / (2\varepsilon^2)$
  ($\varepsilon=\delta=0.05$) and complexity penalty
  $\sqrt{(-\log\pi(k^*)+\log(2\sqrt{N}/\delta))/(2N)}$
  at $N=300$ for Catalan ($\gamma=1$) vs Chipman prior.}
\label{tab:pacbayes}
\begin{tabular}{lccccc}
\toprule
& \multicolumn{2}{c}{$\Nmin$} & \multicolumn{2}{c}{Penalty at $N=300$} & Advantage \\
\cmidrule(lr){2-3}\cmidrule(lr){4-5}
$k^*$ & Cat $\gamma=1$ & Chipman & Cat $\gamma=1$ & Chipman & \\
\midrule
1 & \textbf{74}  & 599 & \textbf{0.107} & 0.126 & Cat 8.1$\times$ \\
2 & 274 & \textbf{119} & 0.115 & \textbf{0.109} & Chip 2.3$\times$ \\
3 & 613 & \textbf{258} & 0.127 & \textbf{0.114} & Chip 2.4$\times$ \\
4 & 996 & \textbf{478} & 0.138 & \textbf{0.122} & Chip 2.1$\times$ \\
\bottomrule
\end{tabular}
\end{table}

\begin{remark}[Medical imaging regime]
For the $N=40$ knee OA dataset, the relevant oracle is $k^*=1$ or $k^*=2$
(single-leaf or two-leaf tree for a 40-sample dataset with binary outcome).
At $k^*=1$, the Catalan prior's $\Nmin=74$ is within one doubling of the
available data, while Chipman's $\Nmin=599$ is $15\times$ beyond reach.
This provides theoretical justification for the Catalan prior's empirical
advantage in the small-$N$ medical imaging setting.
\end{remark}

% ============================================================
\section{Numerical Studies}
\label{sec:experiments}
% ============================================================

All experiments are implemented in Python using NumPy and SciPy;
code and pre-computed results are publicly available at \url{https://github.com/vitsch/jbdt}.
Random seeds are fixed ($\texttt{seed}=42$).

\subsection{Experiment A: WAIC-LOO gap}

We simulate single-leaf DM observations with $C=2$, $\alpha_0=1$,
$p_{\rm true}=(0.5,0.5)$, varying $N \in \{10,20,40,80,160,320,640,1280\}$,
with 500 replicates.
Figure~\ref{fig:waic_gap} (top) shows the total gap
$\waic_{\rm dev} - \loo_{\rm dev}$ converging to $-4$ as $N\to\infty$,
closely tracking the formula $-4N/(N+\alpha_0)$ (maximum deviation $0.03$).
The per-observation gap (bottom) follows the $1/N$ decay.

\subsection{Experiment B: Posterior concentration}

Table~\ref{tab:concentration} summarises $\Pi_N(k^*)$ for $k^*=3$, $C=2$
with informative leaves $[0.1/0.9,\,0.5/0.5,\,0.9/0.1]$, averaged over 50 datasets.
The crossover from Chipman to Catalan advantage occurs at $N\approx 126$
(Catalan $\gamma=1$), consistent with Theorem~\ref{thm:posterior_concentration}.
The universal constant $\reff/(e^{-\gamma}/4) = 1.228$ (Proposition~\ref{prop:catalan_tail}c)
is confirmed across $\gamma \in [0.25, 3.0]$ to $6$ decimal places.

\subsection{Experiment C: BMA oracle inequality}

Table~\ref{tab:bma_oracle} confirms Theorem~\ref{thm:oracle}:
(a) $J_{\rm bound} > \Rbma$ in all 80-replicate experiments (ratios 1.06--1.21);
(b) $\Delta_{\rm Jensen}$ exceeds the excess risk in all cases (ratios 1.12--1.31);
(c) excess risk falls as $c/N$ (56$\times$ reduction from $N=20$ to $N=800$,
vs $40\times$ increase in $N$).

\subsection{Experiment D: WAIC model selection}

With $k^*=3$, the WAIC selection rate (probability that $k^*$ has minimum WAIC)
reaches $100\%$ at $N=500$, matching the theoretical prediction
$N \geq \Nmin(k^*=3) \approx 5.41/\Delta$ with $\Delta \approx 0.01$.

\subsection{Experiment E: PAC-Bayes penalties}

Figure~\ref{fig:pacbayes} (omitted; available in the supplementary code)
plots the PAC-Bayes complexity penalty versus $N$ for $k^*=1,2,3$
under Catalan and Chipman.
At all $N$, the Catalan penalty for $k^*=1$ is $14$--$17\%$ lower than Chipman,
while Chipman is $10\%$ lower for $k^*=3$---consistent with the $\Nmin$ ratios
in Table~\ref{tab:pacbayes}.

% ============================================================
\section{Discussion}
\label{sec:discussion}
% ============================================================

We have established four theorems connecting the DM-leaf likelihood,
Catalan-exponential prior, WAIC model selection, and PAC-Bayes
generalisation bounds for JBDT.

\paragraph{Practical implications.}
The $\Nmin \approx 5.41/\Delta$ formula provides a concrete design criterion
for JBDT practitioners: collect at least $5/\Delta$ observations before
relying on WAIC-weighted BMA.
For the knee OA setting with $\Delta \approx 0.13$, this gives $N_{\min} \approx 40$.
Collecting $N=100$ or more observations per joint compartment would make BMA
fully reliable.

\paragraph{Prior choice.}
The Catalan prior is preferable when the true model is sparse ($k^*=1,2$),
which is the most common scenario for small-$N$ medical imaging.
For denser true models ($k^*\geq 3$) with larger datasets, the Chipman prior
may be more appropriate from a PAC-Bayes standpoint.
A natural extension is a hierarchical prior that learns the effective
decay rate $\gamma$ from data.

\paragraph{Extensions.}
The WAIC gap theorem (Proposition~\ref{prop:gap}) extends immediately to
multi-leaf trees with non-overlapping leaves, since the per-leaf contributions
are independent given the partition.
The oracle inequality (Theorem~\ref{thm:oracle}) extends to BMA over Zernike
feature orders \citep{jakaite2025} by treating each Zernike configuration
as a model family.
The decision-theoretic analysis of Supplementary Appendix~C connects
these results to asymmetric-loss classification, showing that DM-calibrated
posteriors enable an $8.1\times$ reduction in expected decision cost for OA
detection at the optimal threshold $\tau^* = C_{FP}/(C_{FP}+C_{FN}) = 1/3$.

\paragraph{Limitations.}
Our results are exact for the DM leaf model; extension to continuous leaf
models (e.g.\ Gaussian) would require different variance calculations in the WAIC gap.
The PAC-Bayes bound \eqref{eq:pacbayes} is known to be vacuous for large models;
for JBDT with $k \leq 5$ leaves and $N = 40$, the penalty term is $0.27$--$0.45$,
which is non-trivial but smaller than the $0.5$ random-guessing baseline.

% ============================================================
\bibliographystyle{ba}
\bibliography{references}

\end{document}
